% Generated draft: requires mathematical and source review.
\documentclass[11pt]{article}
\usepackage[margin=25mm]{geometry}
\usepackage{fontspec}
\setmainfont{Latin Modern Roman}
\usepackage{amsmath,amssymb,mathtools,unicode-math}
\setmathfont{Latin Modern Math}
\usepackage{enumitem,needspace,xurl,hyperref}
\setlength{\parindent}{0pt}
\setlength{\parskip}{5pt}
\setlength{\emergencystretch}{3em}
\newcommand{\N}{\mathbb{N}}
\newcommand{\Z}{\mathbb{Z}}
\newcommand{\Q}{\mathbb{Q}}
\newcommand{\R}{\mathbb{R}}
\newcommand{\C}{\mathbb{C}}
\newcommand{\F}{\mathbb{F}}
\newcommand{\A}{\mathbb{A}}
\newcommand{\PP}{\mathbb{P}}
\newcommand{\E}{\mathbb{E}}
\newcommand{\eps}{\varepsilon}
\newcommand{\dd}{\,\mathrm{d}}
\begin{document}
\section*{UnsolvedMath ID: 30002435}
\textbf{Source problem code:} catalogue entry
\subsection{Definitions and mathematical statement}
\begin{enumerate}
\item Let \(p\) be a prime, let \(\mathbb F_p\) be the field with \(p\) elements, and write \(\mathbb F_p^*=\mathbb F_p\setminus\{0\}\) for its multiplicative group. For arbitrary polynomials \(f,g\in\mathbb F_p[X]\) and multiplicative subgroups \(G,H\leq\mathbb F_p^*\), define
\[
f(G)=\{f(u):u\in G\},\qquad g(H)=\{g(v):v\in H\},
\]
and
\[
I_p(f,g;G,H)=\#\bigl(f(G)\cap g(H)\bigr).
\]
The first question asks for estimates for \(I_p(f,g;G,H)\). This counts distinct common image values, without multiplicity. The supplied statement imposes no nonconstancy, degree, subgroup-size, or nondegeneracy hypotheses. It specifies neither a proposed inequality nor a desired range, sharpness, or uniformity for the estimates, so a more specific conjectural bound cannot be extracted from the record.

\item Let \(q\) be a prime power, let \(T\) be an indeterminate, and put \(K=\mathbb F_q(T)\). Let \(F\in K[X,Y]\) be an absolutely irreducible polynomial of degree \(d\); absolute irreducibility means irreducibility in \(\overline K[X,Y]\), where \(\overline K\) is an algebraic closure of \(K\). Here degree is interpreted as total degree in \(X,Y\); the supplied record does not explicitly specify this convention. For an integer \(n\geq0\), define
\[
N_F(n)=\#\left\{(x,y)\in\mathbb F_q[T]^2:
\deg_T x\leq n,\ \deg_T y\leq n,\ F(x,y)=0\text{ in }K\right\}.
\]
Use the convention \(\deg_T0=-\infty\), so that zero coordinates are included. The second question asks for upper bounds for \(N_F(n)\) as \(n\to\infty\). For this counting function, \(q\) and \(F\) are fixed while \(n\) varies. The record supplies no target growth rate and does not state whether the requested bounds must be uniform in \(q\), in \(F\) of degree \(d\), or in the coefficients of \(F\), nor how constants may depend on these data. No additional hypotheses on the curve are supplied.
\end{enumerate}
These are two separate estimation questions. The supplied registry labels the record open and its statement \texttt{corrected\_verified}; these are supplied curation labels, not independent mathematical verification.
\subsection{Short English statement}
The first question asks how many distinct values two polynomials can share when evaluated on two multiplicative subgroups of a prime field. The second asks how many ordered pairs of polynomials over \(\mathbb F_q\), each of degree at most \(n\), lie on a fixed absolutely irreducible curve over \(\mathbb F_q(T)\), as \(n\) grows. Neither question includes a specific proposed bound or precise uniformity requirements.
\subsection{Sources}
\begin{itemize}
\item Local UnsolvedMath registry, \texttt{data/unsolvedmath/problems.json}, permanent ID \(30002435\), source problem code \texttt{catalogue entry}, title ``Polynomial Images of Subgroups and Function-Field Curve Points,'' in the Oberwolfach Reports Open Problems collection. This draft uses the supplied complete source record; the local file was not independently inspected.
\item Supplied report citation: \emph{Analytic Number Theory (2014)}, Oberwolfach Reports, pp.~2963--3037, DOI \texttt{10.4171/owr/2013/51}, \url{https://doi.org/10.4171/owr/2013/51}. The supplied literature-source label is ``original OWR report.'' The background describes this collection as drawn from workshop reports and their open-problem discussions. No individual proposer is identified. The registry gives proposed year \(2013\), while its report citation uses \(2014\); both are retained without reconciliation.
\item Supplied literature source labelled ``Shparlinski polynomial values'': \url{https://arxiv.org/abs/1401.0964}. No fuller bibliographic citation or precise supporting result is supplied.
\item Supplied literature assessment, checked \texttt{2026-08-21}: special cases have bounds, while the general two-polynomial, two-subgroup intersection problem remains unresolved. This is supplied dated curation derived from the OWR curation, not independently verified. The supplied background explicitly calls for independent verification of time-sensitive status and cited claims. It gives no separate literature assessment for the function-field counting question.
\item The registry states that the clean statement was checked against the official report and restores its hypotheses and notation. This is a supplied statement-verification claim; no source was visited or independently checked in preparing this draft.
\end{itemize}
\end{document}
